% Source-derived chapter 01: Introduction
% Original source: intersection-complexes-unramified-l-factors
\chapter{Introduction}
\label{intersection-complexes-unramified-l-factors-chapter-01}
\source{intersection-complexes-unramified-l-factors}

\begin{remark}[Source section]
\label{intersection-complexes-unramified-l-factors-chapter-01-source}
\source{intersection-complexes-unramified-l-factors}
\source{intersection-complexes-unramified-l-factors}
This chapter reproduces the corresponding section of the retrieved
LaTeX source, including its definitions, statements, examples, and
proofs. It has not yet been translated into Lean.
\notready
\end{remark}

% The source section title was: Introduction
\subsection{Arc spaces and their IC functions}


Let $\mathbb F$ be a finite field, $G$ a connected reductive group over $\mathbb F$, and $X$ an affine spherical $G$-variety. The formal arc space $\msf L^+X$ is the infinite-dimensional scheme that represents the functor $R\mapsto X(R\tbrac t)$, and it is singular, if $X$ is. However, its singularities at ``generic'' $\mathbb F$-points (namely, arcs $\Spec \mathbb F\tbrac t \to X$ which generically lie in the smooth locus $X^{\text{sm}}$) are of finite type, according to the theorem of Grinberg--Kazhdan and Drinfeld, and this allows one to define an \emph{IC function} \cite{BNS}, that is, the function that should correspond under Frobenius trace to the ``intersection complex'' of $\msf L^+X$. This is a function $\Phi_0$ on $X(\mathfrak o)\cap X^{\text{sm}}(F)$ (where $\mathfrak o=\mathbb F\tbrac t$ and $F = \mathbb F\lbrac t$), and it was conjectured in \cite{SaRS} (before the rigorous definition of this function was available) that it is related to special values of $L$-functions. Such a relation was established in \cite{BNS} for toric varieties and certain group embeddings, termed $L$-monoids, generalizing the local unramified Godement--Jacquet theory. One goal of the present paper is to prove such relations, for a different and broad class of spherical varieties. 

In order to compute the IC function, we need to work with finite-dimensional global models of the arc space, or rather, the arc space over the algebraic closure $k$ of $\mathbb F$; from now on, we  base change to $k$,  without changing notation. However, for motivational purposes, let us here pretend that $\msf L^+X$ were already finite-dimensional, and the intersection complex on it were defined as a $\ol\bbQ_\ell$-valued constructible (derived) sheaf, for some prime $\ell$ different from the characteristic of $\mathbb F$. We would normalize such a sheaf to be constant in degree zero over the arc space of the smooth locus $\msf L^+X^{\text{sm}}$, and we would normalize the intersection complex of a substratum $S\subset \msf L^+X$ of codimension $d$ to be $\ol\bbQ_\ell(\frac{d}{2})[d]$ on its smooth locus, where $\ol\bbQ_\ell(\frac{1}{2})$ denotes a chosen square root of the cyclotomic Tate twist. 

The IC function is computed via some version of a ``Satake transform'' for the spherical variety --- we will recover the original function from its Satake transform in Corollary \ref{cor:asymptotics}. Let $B\supset N$ be a Borel subgroup of $G$, and its unipotent radical, and consider the invariant-theoretic quotient $X\sslash N = \Spec k[X]^N$, which is an affine embedding of a quotient $T_X$ of the Cartan $T=B/N$. In this paper, we restrict ourselves to varieties where $B$ acts freely on an open subset $X^\circ$ of $X$, so let us already make this assumption for notational simplicity. Then $T_X=T$, and we fix a base point to identify $T$ as the open orbit in $X\sslash N$. In fact, our assumption on $X$ is stronger, requiring that the (Gaitsgory--Nadler) \emph{dual group} of $X$, $\check G_X$, is equal to the Langlands dual group of $G$ --- this condition is equivalent\footnote{To be precise, we are referring to the modification of the Gaitsgory--Nadler dual group described in \cite{SV}, because the Gaitsgory--Nadler dual group would be $\check G$ even if the stabilizers were \emph{normalizers} of tori. This small distinction is important, and such cases (for example, $\mathrm{O}_n\backslash \GL_n$) are \emph{not} expected to be directly related to $L$-functions, and not included in the present paper.} to the following:
\begin{equation}
\label{equation-typeT}
\source{intersection-complexes-unramified-l-factors}
\parbox{0.8\linewidth}{$B$ acts freely on $X^\circ$, and for every simple root $\alpha$, if $P_\alpha$ is the parabolic generated by $B$ and the root space of $-\alpha$, the stabilizer of a point in the open $P_\alpha$-orbit $X^\circ P_\alpha$ is a torus (necessarily one-dimensional).}
\end{equation}

The pushforward map $\pi: X\to X\sslash N$ induces a map between arc spaces. Under our assumptions, the ``generic'' $\msf L^+T$-orbits on the arc space of $X\sslash N$, that is, those corresponding to arcs whose generic fiber lands in the open $T$-orbit, are naturally parametrized by a strictly convex (i.e., not containing non-trivial subgroups) submonoid $\mathfrak c_X\subset \check\Lambda$ of the cocharacter group of $T$, with $\check\lambda\in\mathfrak c_X$ corresponding to the image $t^{\ch\lambda}:=\check\lambda(t)$ of a uniformizer (acting on a fixed base point of $X\sslash N$). The pushforward $\pi_! \IC_{\msf L^+X}$ of the IC sheaf is $\msf L^+T$-equivariant, and under Frobenius trace translates to a $T(\mathfrak o)$-invariant function on $X\sslash N (\mathfrak o) \cap T(F)$, which will be denoted by $\pi_! \Phi_0$. Explicitly, 
\[ \pi_! \Phi_0 (a) = \int_{N(F)} \Phi_0(an) dn,\] 
is a Radon transform on the spherical variety, i.e., the integral of the IC function $\Phi_0$ over \emph{generic horocycles}, that is, over the fibers of the map $X(\mathfrak o)\to X\sslash N(\mathfrak o)$, where the Haar measure on $N(F)$ is so that $dn(N(\mathfrak o))=1$. 

\smallskip

This integral is really a finite sum, hence makes sense over $\ol\bbQ_\ell$, but let us for simplicity choose an isomorphism $\ol\bbQ_\ell \simeq \mathbb C$, such that the geometric Frobenius morphism $\Fr$ acts on the chosen half-Tate twist $\ol\bbQ_\ell(\frac{1}{2})$ by $q^\frac{1}{2}$. 
Our results and conjectures are best expressed under the assumption that $X$ carries a $G$-eigen-volume form;  we will assume this for the rest of the introduction. The absolute value of the volume form is a $G(F)$-eigenmeasure on the open orbit $X^\bullet(F)$, whose  eigencharacter we will denote by $\eta$. (We will not impose this assumption in the rest of the paper, but see Remark \ref{rem:length}.) Then, one should consider the following normalized form of the integral above, analogous to the standard normalization of the Satake isomorphism:
\begin{equation}
(\eta\delta)^\frac{1}{2}(a) \pi_! \Phi_0 (a) = (\eta\delta)^\frac{1}{2}(a)\int_{N(F)} \Phi_0(an) dn,
\end{equation}
where $\delta = |e^{2\rho_G}|$ is the modular character of the Borel subgroup.\footnote{We use additive notation for the character group $\ch\Lambda$ of $\ch T$, so $e^{\ch\nu}$ will denote the actual morphism $\ch T\to \mathbb G_m$ corresponding to $\ch\nu\in \ch\Lambda$.}


Ideally, we would like to prove a conjecture such as the following. To formulate it, recall that $T(\mathfrak o)$-orbits on $T(F)$ are parametrized by Galois-fixed (that is, $\Fr$-fixed) elements of $\ch\Lambda$.

\begin{conj}
\source{intersection-complexes-unramified-l-factors}
\notready
\label{conjecture-pushforward-IC}
\source{intersection-complexes-unramified-l-factors}
There is a symplectic representation $\rho_X$ of the $L$-group, $\rho_X: {^LG}=\check G \rtimes \left<\Fr\right> \to \GL(V_X)$, and a ${^LT}$-stable polarization $V_X=V_X^+\oplus V_X^-$, such that the multiset $\mB^+$ of $\check T$-weights of $V_X^+$ belongs to $\mathfrak c_X$, and the pushforward of the IC function satisfies:
 \begin{equation}
 \label{equation-pushforward-IC}
\source{intersection-complexes-unramified-l-factors}
  (\eta\delta)^\frac{1}{2} \cdot \pi_! \Phi_0 = \left(\tr_{\check T}(\Fr, \Sym^\bullet (\check{\mathfrak n}(1)))\right)^{-1} \cdot \tr_{\check T}(\Fr, \Sym^\bullet (V_X^+)).
 \end{equation}
 Here, for a representation $V$ of ${^LT}=\check T\rtimes\left<\Fr\right>$, the expression $\tr_{\check T}(\Fr,V)$ denotes the function on $\ch\Lambda^{\Fr}$ whose value on $\ch\lambda$ is equal to the trace of geometric Frobenius $\Fr$ on the $(\ch T, \ch\lambda)$-eigenspace of $V$. 
\end{conj}

From the point of view of number theory, the spherical varieties satisfying our assumption $\check G_X=\check G$ give, in some sense, the most interesting periods, because they are associated to $L$-values at the center of the critical strip. Indeed, one should always be able to choose the eigencharacter $\eta$ such that the Frobenius morphism acts on $V_X^+$ by permuting elements of a basis and scaling by $q^\frac{1}{2}$. 
For example, when $G$ is split and the colors (see below) of $X$ are all defined over $\mathbb F$, this permutation action should be trivial, and \eqref{equation-pushforward-IC} should read:
\begin{equation}\label{equation-pushforward-IC-expl}
\source{intersection-complexes-unramified-l-factors}
(\eta\delta)^\frac{1}{2} \pi_! \Phi_0 = \frac{\prod_{\check\alpha\in\check\Phi^+} (1-q^{-1} e^{\check\alpha})}{\prod_{\check\lambda \in \mB^+} (1-q^{-\frac{1}{2}} e^{\check\lambda})},  
\end{equation}
where $\mB^+$ is the multiset of weights, as in the conjecture, and $e^{\ch\lambda}$ denotes the characteristic function of the $T(\mf o)$-orbit of $t^{\ch\lambda}$.

We will explain the relationship of this conjecture to various conjectures of arithmetic and geometric origin below. We do not quite prove the conjecture in all cases, but we determine the weights of $\rho_X$ 
(in terms of $X$) 
and in the cases where $\rho_X$ is minuscule, we prove the conjecture (Corollary~\ref{cor:minuscule}). 
It is helpful to distinguish between the special case when $X = \overline{H\backslash G}^\aff$ is the affine closure of a homogeneous quasiaffine variety, and the general case. In the special case, let us also assume, at first, that 
the monoid $\mf c_X$ is freely generated with a basis $\ch\nu_1,\dotsc, \ch\nu_r$, so 
$X\sslash N$ may be identified with $\mbb A^r$. This condition can always be achieved by passing to an abelian cover of $H\backslash G$ and taking its affine closure, see \S \ref{sect:freemonoid}. In that case, the generators $\ch\nu_i$ are the valuations associated to the \emph{colors}, that is, the prime $B$-stable divisors of $H\backslash G$. 

\begin{thm}[See \S \ref{sect:finitefields}]
\source{intersection-complexes-unramified-l-factors}
\notready
\label{theorem-affine-closure}
\source{intersection-complexes-unramified-l-factors}
 Assume that $X = \overline{H\backslash G}^\aff$ satisfies the conditions above ($T_X=T$, $\check G_X=\check G$ and $\mf c_X \cong \mbb N^r$ is free). Then there is a $(\check T\rtimes\left<\Fr\right>)$-representation $V^+_X$ satisfying:
 \begin{enumerate}
  \item the $\check T$-weights of $V^+_X$ belong to $\mathfrak c_X\sm 0$;
  \item the set of weights of $V^+_X \oplus (V^+_X)^*$ (without multiplicities) equals the set of weights of a $\ch G$-representation $\rho_X$;
  \item the dimensions of the weight spaces of $V^+_X \oplus (V^+_X)^*$ are invariant under the Weyl group $W$ of $G$;
  \item the weight spaces in $V^+_X$ of the basis elements $\check\nu_1,\dots,\check\nu_r$ of $\mathfrak c_X$ have multiplicity one,
  \item under the Frobenius action, we have 
  \begin{equation}\label{e:Frobaction}
\source{intersection-complexes-unramified-l-factors}
   V^+_X= \bigoplus_{\mB^+} \ol\bbQ_\ell (\smallfrac{1}{2}),
\end{equation}
 for some permutation action of $\Fr$ on the multiset $\mB^+$ of weights, compatible with its action on $\ch T$,    
\end{enumerate}
 such that the pushforward $\pi_! \Phi_0$ of the IC function\footnote{Under the assumptions of the theorem, the restriction of the eigencharacter $\eta$ to the colors $\ch\nu_i$ is uniquely determined, see Remark \ref{rem:length}.} satisfies the formula \eqref{equation-pushforward-IC} above.
\end{thm}

In fact, we show more: we endow the multiset $\mB := \mB^+ \sqcup (-\mB^+)$ of weights of $V_X$ with the structure of a \emph{Kashiwara crystal}, see Theorem \ref{theorem-crystal} below, and show that, if Conjecture~\ref{conjecture-pushforward-IC} holds (equivalently, if the crystal is the one corresponding to the crystal basis of a finite-dimensional $\ch G$-module),
then $\rho_X$ must be the direct sum of the irreducible $\ch G$-modules with highest weights 
in $\ch\Lambda^+ \cap W \{ \ch\nu_1,\dotsc,\ch\nu_r \}$, each with multiplicity one 
(see Remark~\ref{rem:conj}).
In other words, the highest weights of $\rho_X$ are the dominant coweights that are Weyl translates
of the basis elements $\ch\nu_1,\dotsc,\ch\nu_r$. 

\smallskip

As we already mentioned, Theorem~\ref{theorem-affine-closure} implies Conjecture~\ref{conjecture-pushforward-IC} when $\rho_X$ is minuscule. 
In particular, when $H\bs G$ is itself affine (equivalently, $H$ is reductive: see \cite{Lun73,Richardson}),
we observe that $\rho_X$ is always minuscule (Corollary~\ref{cor:Hreductive}). In this case
Conjecture~\ref{conjecture-pushforward-IC} was previously proved by \cite[Theorem 7.2.1]{SaSph}, under some additional assumptions,
and Theorem~\ref{theorem-affine-closure} gives a geometric interpretation of this result.


For an example when $H\bs G$ is not affine:

\begin{eg}
\source{intersection-complexes-unramified-l-factors}
\notready\label{example-nfold}
\source{intersection-complexes-unramified-l-factors}
Let $X^\bullet=$ the quotient of $\SL_2^n$ by the unipotent subgroup 
\[ H_0= \left\{ \left.\left(\begin{array}{cc} 1 \\ x_1 & 1\end{array}\right) \times \left(\begin{array}{cc} 1 \\ x_2 & 1\end{array}\right) \times \cdots \times\left(\begin{array}{cc} 1 \\ x_n & 1\end{array}\right) \right| x_1+ x_2 + \dots+ x_n = 0\right\},\]
under the action of $G = $ the quotient of $\mathbb G_m \times \SL_2^n$ by the diagonal copy of $\mu_2$, where $a\in \mathbb G_m$ acts as left multiplication by $\begin{pmatrix} a^{-1} \\ & a\end{pmatrix}$. 
Let $X$ be the affine closure of $X^\bullet$. Denoting by $\check m$ the identity cocharacter of $\mathbb G_m$, the monoid $\mathfrak c_X$ is freely generated by the coweights
\[ \frac{\check\alpha_1+\check\alpha_2+\dots+\check\alpha_n - \check m}{2}\]
and 
\[ \frac{-\check\alpha_1- \dots - \check\alpha_{i-1} + \check\alpha_i - \dots -\check\alpha_n + \check m}{2}, \,\, i =1,\dots, n,\]
see Remark \ref{rem:valuation}. These are minuscule weights of $\check G= \GL_2\times_{\det} \GL_2 \times_{\det} \cdots \times_{\det}\GL_2$, and in that case Conjecture \ref{conjecture-pushforward-IC} holds, confirming an expectation of \cite[\S 4.5]{SaRS}.  
\end{eg}

For the general case, $X$ contains an open $G$-orbit $X^\bullet = H\backslash G$, which we will assume to satisfy the conditions of Theorem \ref{theorem-affine-closure}, except perhaps for the freeness of $\mathfrak c_{X^\bullet}$. In that case, the free monoid $\mathbb N^{\mathcal D}$, where $\mathcal D$ denotes the set of colors, maps through the valuation map to $\mathfrak c_{X^\bullet}$, and $\mathfrak c_X$ is generated by its image and a minimal set $\Cal D^G_{\mathrm{sat}}(X)=\{\check \theta_1, \dots,\check\theta_d\}$ of distinct \emph{antidominant} elements of the cocharacter group $\check\Lambda$ of $T$. For each $\check\theta_i$, we let $V^{\check\theta_i}$ be the irreducible module of $\check G$ with lowest weight $\check\theta_i$, and assume that the eigencharacter $\eta$ of the $G$-eigenmeasure on $X(F)$ is of the form $|e^{\mathfrak h}|$ for some algebraic character $\mathfrak h \in \Lambda^W$. 

\begin{thm}[See \S \ref{sect:finitefields}]
\source{intersection-complexes-unramified-l-factors}
\notready
\label{theorem-general}
\source{intersection-complexes-unramified-l-factors}
In the setting above, if $V_{X^\bullet}^{+}$ denotes the $\check T$-representation for $\overline{X^\bullet}^\aff$ described in Theorem \ref{theorem-affine-closure},\footnote{See \S \ref{sect:freemonoid} for a reduction to the case where $\mathfrak c_{X^\bullet}$ is free.} then the pushforward $  (\eta\delta)^\frac{1}{2}\cdot \pi_! \Phi_0 $ of the IC function for $\msf L^+X$ is given by \eqref{equation-pushforward-IC} with 
\begin{equation}
 V^+_X = V_{X^\bullet}^{+} \oplus \bigoplus_i V^{\check\theta_i}\left(\frac{\left<\mathfrak h + 2\rho_G,\ch\theta_i\right>}{2}\right).
\end{equation}
\end{thm}

Notice that the set $\Cal D^G_{\mathrm{sat}}(X)$ is stable under the Frobenius morphism. The action of Frobenius on the sum of $V^{\check\theta_i}$'s is the one obtained by identifying the crystal basis of this space with a set of subvarieties of the affine Grassmannian (see Section \ref{sect:crystal}), and considering the Frobenius action on those.

The reader should compare the passage from Theorem \ref{theorem-affine-closure} to Theorem \ref{theorem-general} to the passage from a reductive group $G$ to an $L$-monoid $X$ in \cite[Theorem 4.1]{BNS, BNS-erratum}. While the result is similar, however, the straightforward proof of \cite{BNS} uses the monoid structure on $X$ in a crucial way, and cannot be used here. 

\medskip

In \S \ref{subsection-introduction-global} below we will describe the sheaf-theoretic statements of these theorems, in the setting of appropriate finite type models, the \emph{Zastava spaces} for $X$ and $X\sslash N$. Before we do that, let us relate the results above to conjectures in number theory and geometry.


\subsection{Asymptotics and $L$-functions}

The Radon transform $\pi_! \Phi_0$ of the IC function (also known as ``basic function'') under the map $X\to X\sslash N$ admits various interpretations in terms of harmonic analysis, and, in particular,  allows us to compute the Plancherel density of the basic function,
\begin{equation}\label{Plancherel} \Vert \Phi_0 \Vert^2 = \int_{\ch T^1/W} \Omega(\chi) d\chi,
\source{intersection-complexes-unramified-l-factors}
\end{equation}
that is, the decomposition of its norm in the space $L^2(X^\bullet(F))$ (with respect to the fixed eigenmeasure) in terms of seminorms $\Vert \bullet \Vert_\chi$ that factor through eigenquotients for the action of the unramified Hecke algebra. The variable $\chi\in \ch T^1$ here denotes an unramified unitary character of $T(F)$ (an element of the maximal compact subgroup of the complex dual Cartan), which modulo the action of $W$ represents the Satake parameter of such an eigenquotient, and $\Omega(\chi) = \Vert \Phi_0\Vert^2_\chi$. 

The passage from $\pi_! \Phi_0$ to the Plancherel formula \eqref{Plancherel} is achieved through the theory of \emph{asymptotics}, which is of geometric interest because of its relation to \emph{nearby cycles}. For an affine spherical variety $X$ over a field $k$ in characteristic zero, one can define its \emph{horospherical degeneration} (or asymptotic cone) $X_\emptyset$, by passing to the associated graded of the coordinate ring $k[X]$ as a $G$-module. A similar degeneration exists under some assumptions in positive characteristic, see \S\ref{sect:degeneration}. Under our current assumptions, its open $G$-orbit $X_\emptyset^\bullet$ is isomorphic to $N^-\backslash G$, where we use $N^-$ to denote the unipotent radical of a Borel $B^-$ opposite to $B$ --- an expository choice 
without mathematical significance, which we will use to identify the abstract Cartan $T=B/N$ as an automorphism group of $X_\emptyset$ by identifying it with the torus $B^-\cap B$ (the latter acting ``on the left'' on $N^-\backslash G$). 

The theory of asymptotics states that there is a canonical morphism 
\[ e_\emptyset^*: C^\infty(X^\bullet(F)) \to C^\infty(X_\emptyset^\bullet(F))\] 
which describes the behavior of any function ``at infinity'', see \cite[\S 5]{SV}. Of interest to us is that the spaces $X\sslash N$ and $X_\emptyset\sslash N$ are \emph{canonically} identified, and the corresponding pushforwards $\pi_!$ and $\pi_{\emptyset !}$ satisfy
\begin{equation}\label{Radon-asymptotics} \pi_! = \pi_{\emptyset !} \circ e_\emptyset^*
\source{intersection-complexes-unramified-l-factors}
\end{equation}
 (for appropriate functions, in order to ensure convergence), see 
\cite[Proposition 5.4.6]{SV}. 

Inverting the Radon transform, \eqref{Radon-asymptotics} leads to a calculation of the asymptotics of the basic function. This, in turn, leads to a calculation of the basic function itself, as a function on $X^\bullet(F)/G(\mf o)$. To express both, we note that there is a natural parametrization of the coset space $X_\emptyset^\bullet(F)/G(\mf o)$ by the Frobenius-stable elements of the coweight lattice $\ch\Lambda$ (simply assigning $\ch\theta \in \ch\Lambda^\Fr$ to the orbit of $N^-t^{\ch\theta}$), and of the coset space $X^\bullet(F)/G(\mf o)$ by the Frobenius-stable elements of its antidominant submonoid $\ch\Lambda^-$ --- this is explained in Theorem \ref{thm:Gorbits}, in the geometric setting, from which we will borrow the notation $x_0 t^{\ch\theta}$ to represent the orbit corresponding to $\ch\theta\in \ch\Lambda^{-,\Fr}$. Then we have:



\begin{cor}[See \S \ref{sect:asymptotics}]
\source{intersection-complexes-unramified-l-factors}
\notready \label{cor:asymptotics}
\source{intersection-complexes-unramified-l-factors}
In the setting of Theorem \ref{theorem-general}, we have 
 \begin{equation}
 \label{equation-asymptotics-IC}
\source{intersection-complexes-unramified-l-factors}
(\eta\delta)^\frac{1}{2}(a)  e_\emptyset^* \Phi_0 (N^- a G(\mf o)) = \tr_{\check T}(\Fr, \Sym^\bullet (\check{\mathfrak n}))^{-1} \cdot \tr_{\check T}(\Fr, \Sym^\bullet (V_X^+)) ,
\end{equation}
and 
\begin{equation}
 \label{equation-basicfunction}
\source{intersection-complexes-unramified-l-factors}
(\eta\delta)^\frac{1}{2}(a)  \Phi_0 (x_0 a G(\mf o)) = \left. \tr_{\check T}(\Fr, \Sym^\bullet (\check{\mathfrak n}))^{-1} \cdot \tr_{\check T}(\Fr, \Sym^\bullet (V_X^+))\right|_{\ch\Lambda^{-,\Fr}},
\end{equation}
where we use the notation $\tr_{\ch T}(\dotsb)$, as in Conjecture \ref{conjecture-pushforward-IC}, to represent a $T(\mf o)$-invariant function on $T(F)$. 
\end{cor}


For example, in the setting of  \eqref{equation-pushforward-IC-expl}, the right hand side reads
\[\frac{\prod_{\check\alpha\in\check\Phi^+} (1- e^{\check\alpha})}{\prod_{\check\lambda \in \mB^+} (1-q^{-\frac{1}{2}} e^{\check\lambda})}.\]



Finally, the Plancherel decomposition of $\Phi_0$ coincides\footnote{See the proof of Proposition \ref{prop:localzeta}.} up to the factor $|W|^{-1}$ with that of $e_\emptyset^*\Phi_0$, which by Mellin transform on $X_\emptyset^\bullet$ (with respect to the action of $T$) gives:
\begin{equation}
 \label{equation-Plancherel-IC}
\source{intersection-complexes-unramified-l-factors}
\Vert \Phi_0\Vert^2 = \int_{\hat T/W} \frac{\prod_{\check\alpha\in\check\Phi} (1- e^{\check\alpha}(\chi))}{\prod_{\check\lambda\in S} (1-q^{-\frac{1}{2}} e^{\check\lambda}(\chi))} d\chi = \int_{\hat T/W} L(\chi, V_X, \frac{1}{2}) \frac{d\chi}{L(\chi, \check{\mathfrak g}/\check{\mathfrak t},0)}.
\end{equation}

Here $\mB = \mB^+ \sqcup (-\mB^+)$, and  $L(\chi, V_X, 0)$ denotes a local unramified $L$-factor, while the density $\frac{d\chi}{L(\chi, \check{\mathfrak g}/\check{\mathfrak t},0)}$ is the unramified Plancherel measure for $G$. 

The importance of \eqref{equation-Plancherel-IC} for arithmetic is that, according to the generalized Ichino--Ikeda conjecture of \cite{SV}, the quotient of the Plancherel density of $\Phi_0$ by the Plancherel measure of $G$ is related to the local Euler factor of the ``$X$-period integral'' of automorphic forms. A provable case of such an Euler factorization is related to Example \ref{example-nfold}:

\begin{eg}
\source{intersection-complexes-unramified-l-factors}
\notready\label{example-nfold-global}
\source{intersection-complexes-unramified-l-factors}
 Let $G$, $X$ be as in Example \ref{example-nfold}, over the function field $\Bbbk = \mathbb F(C)$ of a curve $C$, and let $\Phi = \prod_v \Phi_v$ be a smooth, factorizable function of moderate growth on the adelic points of the open $G$-orbit $X^\bullet$, such that the support of $\Phi_v$ has compact closure in $X(F_v)$, and for almost every $v$ the function $\Phi_v$ is equal to the IC function of $X$. Let $\Theta_\Phi(g) = \sum_{\gamma\in X^\bullet(\Bbbk)} \Phi(\gamma g)$ be the corresponding theta series, a function on the adelic points of $G$. 
 
 Let $\phi$ be a cusp form, belonging to a cuspidal automorphic representation $\pi =  \bigotimes_v' \pi_v$, and let $W_\phi(g) = \prod_v W_{\phi, v}(g_v) = \int_{N^-(\Bbbk)\backslash N^-(\mbb A)} \phi(ng) \psi(n) dn$ be the Whittaker function of $\phi$ (where $N^-$ is the lower triangular subgroup), with a chosen Euler factorization with $W_{\phi,v}(1)=1$ for almost all $v$. It can be proven by the usual unfolding argument that the pairing
 \[ \int_{G(\Bbbk)\backslash G(\mbb A)} \phi(g) \Theta_\Phi(g) dg\]
 is convergent if the central character of $\phi$ is ``large'' enough (i.e., its restriction to $\mathbb G_m$ has absolute value $|\bullet|^s$ for some $s \gg 0$), and equal to the Euler product of zeta integrals 
 \[ \int_{H_0\backslash G (\Bbbk_v)} W_{\phi,v}(g_v) \Phi_v(g_v) dg_v.\]
 
 Our Theorem \ref{theorem-affine-closure} implies that almost every Euler factor is equal to the local unramified $L$-factor $L(\pi_v, \otimes, 1 - \frac{n}{2})$, where $\otimes$ is the tensor product representation of $\check G$.  We explain this in \S \ref{sect:nfold-explanation}. 
\end{eg}







\medskip

Such global applications are beyond the main focus of this paper. Of more immediate interest here is the relation of the asymptotics map $e_\emptyset^*$ to nearby cycles: Since $X_\emptyset$ is obtained by degenerating the coordinate ring of $X$, there is an associated $\mathbb G_m\times G$-equivariant Rees family $\mf X\to \mathbb A^1$ (depending, really, on the choice of a strictly dominant cocharacter into $T$), whose general fiber is isomorphic to $X$, and whose special fiber is isomorphic to $X_\emptyset$. This also induces a family of arc spaces, or loop spaces $\msf L_{\mbb A^1}\mf X\to \mathbb A^1$ where $\msf L_{\mbb A^1}\mf X$ denotes the family of \emph{fiberwise} loop spaces, not the loop space of $\mf X$. In the context of an appropriate sheaf theory, to be denoted by $D$, this would give rise to a nearby cycles map:
\[ \Psi: D(\msf LX) \to D(\msf LX_\emptyset),\]
whose Frobenius trace is expected to recover the asymptotics map $e_\emptyset^*$. 

After replacing $\msf L^+X, \msf L^+ X_\emptyset, \msf L^+(X\sslash N)$ by finite type models $\sY, \sY_\emptyset, \sA$, respectively (see \S\ref{subsection-introduction-global} below), we prove: 

\begin{thm}[see Theorem~\ref{thm:PsiRadon}]
\source{intersection-complexes-unramified-l-factors}
\notready
\label{thm:nearby-asymptotics}
\source{intersection-complexes-unramified-l-factors}
Let $X$ be an affine spherical variety such that $B$ acts freely on $X^\circ$. 
Then the following triangle of functors commutes up to natural isomorphism:
\[ \begin{tikzcd} 
\rmD^b_c(\sY) \ar[r, "\Psi"] \ar[rd,swap, "\pi_!"] & \rmD^b_c(\sY_\emptyset) \ar[d, "\pi_{\emptyset !}"]\\
& \rmD^b_c(\sA) 
\end{tikzcd}
\]
\end{thm}

The function-theoretic asymptotics map $e_\emptyset^*$ satisfies the same 
commutative triangle \eqref{Radon-asymptotics}, which suggests that nearby cycles $\Psi$
is in a suitable sense the geometrization of the asymptotics map.
This resembles an analogous result \cite[Corollary 6.2]{BFO} in the setting 
of character sheaves.
In the case where $X$ is a reductive group, the nearby cycles of the IC complex 
of (finite type models of) $\msf L^+X$ 
has been computed by S.~Schieder \cite{Sch-SL2,Sch-G,Sch}. 

In Theorem~\ref{thm:nearby}, we compute the intersection complex of the global model $\sM_X$ of the arc space, for the spherical varieties in \eqref{equation-typeT}, at the level of Grothendieck groups. We do this by relating the nearby cycles complex to $\pi_!\IC_\sY$, in a way that corresponds to the known relation \eqref{Radon-asymptotics} between asymptotics and Radon transforms. 
\medskip

Finally, we explain how a formula like \eqref{equation-Plancherel-IC} relates to recent conjectures of Ben-Zvi--Sakellaridis--Venkatesh \cite{BZSV}: According to those conjectures, the formula should follow by applying Frobenius traces on the \emph{endomorphism ring} $\End(\IC_{\msf L^+X})$, where the endomorphism ring is taken in the derived sense, in the dg-category of derived constructible sheaves on $\msf LX/\msf L^+G$. (The proper definition of the intersection complex on the arc space remains conjectural, in the singular setting.) This conjecture identifies (ignoring cohomological grading) 
\begin{equation}\label{equation-BZSV}
\source{intersection-complexes-unramified-l-factors}
  \End( \IC_{\msf L^+X_k} ) \simeq \ol\bbQ_\ell [V_X]^{\check G}
\end{equation}
for some symplectic representation $V_X$ as above. We hope that our results can be related to this conjecture in both directions: Namely, that by relating our results to such endomorphism rings, one can upgrade the $\ch T$-structure of Theorem \ref{theorem-affine-closure} to a $\ch G$-structure, proving Conjecture \ref{conjecture-pushforward-IC}; and, vice versa, one can make progress towards the conjecture of \cite{BZSV} by utilizing our results.

The conjectures of \cite{BZSV}, and the problems addressed by the present paper, can be formulated for more general affine spherical varieties, without the assumption that $\ch G_X = \ch G$. As already mentioned, from the point of view of number theory, the latter case is perhaps the most interesting one, as it corresponds to central values of $L$-functions. In the general case, the vector space $V_X$ appearing in the conjectural relation \eqref{equation-BZSV} needs to be replaced by a Hamiltonian manifold living over the quotient $\ch G_X\backslash \ch G$.


\subsection{Zastava spaces and the main theorems in terms of sheaves}
\label{subsection-introduction-global}
\source{intersection-complexes-unramified-l-factors}

From now on, we work over the algebraic closure $k$ of the finite field $\mathbb F$, or over an algebraically closed field $k$ in characteristic zero. When $X$ is defined over a finite field $\mathbb F$, we will keep track of Weil structures on our sheaves, which will always have the form of half-integral Tate twists, where, as mentioned, $(\frac{1}{2})$ denotes a fixed square root of the cyclotomic twist. The intersection complex of a $d$-dimensional scheme over $k$ will be understood to have stalks $\ol\bbQ_\ell (\frac{d}{2})[d]$ over the smooth locus. 

In order to replace the arc space by a model of finite type, we fix a smooth projective curve $C$ over $k$ (or $\mathbb F$). For an algebraic stack $\Scr X$ and an open substack $\Scr X^\circ \subset \Scr X$,
we will use 
\[ \Maps_\gen(C, \Scr X \supset \Scr X^\circ) \]
to denote the prestack that assigns to a test scheme $S$
the groupoid of maps $C\times S \to \Scr X$
such that the open locus of points sent to $\Scr X^\circ$ 
maps surjectively to $S$. 
Equivalently, these are the maps such that for every geometric point $\bar s \to S$, 
the restricted map $C\times \bar s \to \Scr X$ generically lands in $\Scr X^\circ$.
Since $C$ is smooth, $\Maps_\gen(C, \Scr X \supset \Scr X^\circ)$ is 
an open substack of the prestack $\Maps(C,\Scr X)$. 

\smallskip

Given an affine spherical $G$-variety $X$ with open $G$-orbit $X^\bullet$ and open $B$-orbit $X^\circ$, we consider the following two models for the arc space of $X$:
\begin{itemize}
 \item the Artin stack 
\[ \sM= \sM_{X} = \Maps_\gen(C, X/ G \supset X^\bullet /G), \]
that we will simply refer to as ``the global model'';
 \item the stack
\[ \Scr Y=   \Scr Y_{X} = \Maps_\gen(C, X/B \supset X^\circ/B)  \]
that we will refer to as ``the Zastava model''. In our setting ($X^\circ\cong B$), this turns out to be a scheme. Such a model is often referred to as the ``local model'' for reasons that have to do with factorization structures, but since this can create confusion with the genuinely local arc space, we will avoid such terminology. 
\end{itemize}
For a discussion of why these are indeed formal models of the arc space (in the formal neighborhoods of suitable points), see Theorem~\ref{thm:DGK}, Lemma~\ref{lem:localglobalyoga}.
Note that the choice of a Borel subgroup is immaterial, since $X/B$ can also be written as $(X\times \mathcal B)/G$, where $\mathcal B$ is the flag variety, with $X^\circ/B = (X\times \mathcal B)^\bullet/G$, where $(X\times \mathcal B)^\bullet$ is the open 
orbit under the diagonal $G$-action.  

We will also let $\Scr A$ denote the analog of these models for the toric variety $X\sslash N$, that is,
\[    \Scr A = \Maps_\gen(C, (X\sslash N)/T \supset (X^\circ\sslash N)/T),  \]
and notice that under our assumptions $X^\circ\sslash N\simeq T$. Fixing such an identification, for every $\chi\in \mathfrak c_X^\vee$ 
(the subset of $\Lambda_X$, the character group of $T$, of those elements that are $\ge 0$ on $\mathfrak c_X$), the corresponding map $X\sslash N\to \mathbb G_a$ gives rise to a morphism $\Scr A \to \Maps_\gen(C, \mathbb G_a/\mathbb G_m \supset \mathbb G_m/\mathbb G_m) = \Sym C$, the scheme of effective divisors on the curve. Thus, $\Scr A$ can be thought of as the scheme of $\mathfrak c_X$-valued divisors. This scheme is well understood \cite{BNS}: its normalization is a disjoint union of partially symmetrized powers $C^{\mf P}$ of the curve, indexed by formal $\mathbb N$-linear combinations $\mf P\in \Sym^\infty(\Prim(\mathfrak c_X))$ of the primitive elements of $\mathfrak c_X$ (see \S\ref{def:partition}).

\smallskip

The sheaf-theoretic analog of Theorems \ref{theorem-affine-closure} and \ref{theorem-general} is a statement about the pushforward of the IC sheaf under 
\[ \pi: \Scr Y \to \Scr A.\]
We will only compute $\pi_! \IC_{\Scr Y}$ in the Grothendieck group of sheaves on $\Scr A$ 
(see Corollary~\ref{cor:pibarY}), which is enough to determine the trace of Frobenius on stalks. 
The reason we do not compute the pushforward in the DG category (although this can be done in principle)
is related to the fact that the map $\pi$ is not proper. 
However, one can compactify $\pi$ by considering\footnote{One makes the usual modification to the definition when $[G,G]$ is not simply connected.} the \emph{compactified Zastava space} 
\[ \barY =   \Maps_\gen(C, (\overline{X/N})/T \supset X^\circ/B),  \]
where $\overline{X/N}$ stands for the stack $(X\times \overline{N\backslash G})/G$, where $\overline{N\backslash G}= \Spec k[N\backslash G]$ is the affine closure of $N\backslash G$. 

The difference between $\barY$ and $\Scr Y$ will account for the factor of $\prod_{\check\alpha\in\check\Phi^+} (1-q^{-1} e^{\check\alpha})$ in \eqref{equation-pushforward-IC}. The number theory-minded reader will recognize in this factor, in the case of $G=\SL_2$, the Euler factor of the quotient between Eisenstein series obtained by summing over integral points of $N\backslash \SL_2$, versus integral points of $\overline{N\backslash \SL_2}=\mathbb A^2$. More generally, this is the factor that relates the ``naive'' and ``compactified'' Eisenstein series of \cite{BG}, \cite{BFGM}.

\smallskip

In Proposition \ref{prop:piproper} and Theorem \ref{thm:semismall} we prove:
\begin{thm}
\source{intersection-complexes-unramified-l-factors}
\notready
\label{theorem-semismall}
\source{intersection-complexes-unramified-l-factors}
 The map $\bar\pi: \overline{\Scr Y}\to \Scr A$ is proper and stratified semi-small.
\end{thm}
This is one of the key technical results of this paper, because it allows us to get our hands on the pushforward of the intersection complex, without having a description of the complex itself. 
The assumption that $\ch G_X=\ch G$ is critical for the theorem: the analogous statement 
for the usual Finkelberg--Mirkovi\'c Zastava space (\cite{FM, BFGM}) is far from true. 

The condition of being stratified semi-small is a condition on ``smallness'' of fibers, relative to a fixed stratification which, in this case, is the natural stratification of $\Scr A$ by strata of the form
\[ \iota^{\mf P}: \oo C^{\mf P} \into \Scr A, \] 
where $\mathfrak c_X$-valued divisors take a fixed set of values. (Here, $\oo C^{\mf P}$ denotes the open ``disjoint'' locus in a certain product of symmetric powers of the curve, corresponding to divisors of the form $\sum_{\check\lambda \in\mathfrak c_X} \sum_{i=1}^{N_{\check\lambda}} (v_i) \check\mu$ with all $v_i \in \abs C$ distinct; we will denote by $\bar\iota^{\mf P}$ the natural compactification.)

\smallskip

By the decomposition theorem, stratified semi-smallness ensures that $\bar\pi_! \IC_{\barY}$ is a direct sum of irreducible \emph{perverse} sheaves. By a factorization property of $\barY$, this easily implies an expression for $\bar\pi_! \IC_{\barY}$ of the form
\begin{equation} \label{e:format1-intro}
\source{intersection-complexes-unramified-l-factors}
 \bar\pi_!(\IC_{\barY}) \cong 
\bigoplus_{\mf P} 
\Bigl( \bigotimes_{\ch\lambda} 
\Sym^{N_{\ch\lambda}}(V_{X,\ch\lambda}) \Bigr) \ot 
\bar\iota^{\mf P}_!( \IC_{C^{\mf P}} ),
\end{equation}
(see Proposition \ref{prop:format}), where $V_{X,\ch\lambda}$ are the contributions of \emph{diagonal strata} $\iota^{[\check\lambda]}: C \into \Scr A$, corresponding to divisors supported at one point. Moreover, perversity and an estimate of dimensions imply that these contributions all come from the \emph{top degree} cohomology of those fibers of the map $\Scr Y\to \Scr A$ which (over diagonal strata) are of ``maximal possible dimension'' in terms of the semi-smallness  inequality. (We will discuss this dimension calculus in detail below.) These fibers, over points on the diagonal stratum $\iota^{[\check\lambda]}(C)$ are called the \emph{central fibers}, and denoted by $\msf Y^{\check\lambda}$, with the point not appearing explicitly in the notation. 


\medskip


Therefore, to complete the calculation and prove Conjecture \ref{conjecture-pushforward-IC}, we would need to count, for every $\check\lambda \in \mf c_X$, the irreducible components of the central fiber $\msf Y^{\check\lambda}$ which achieve the maximal dimension, and show that their cardinality is equal to the dimension of the $\check\lambda$-weight space in $V_X^+$, which is ``half'' of a symplectic $\ch G$-representation $\rho_X$. 
Note that $\mf c_X$ is strictly convex, and the weights of $\rho_X$ will have the property that 
they belong to either $\mf c_X\sm 0$ or $-\mf c_X \sm 0$.  Thus, the monoid $\mf c_X$ determines
which ``half'' of $\rho_X$ to consider.

\smallskip

There is a reduction from the general case of Theorem \ref{theorem-general} to the special case of Theorem \ref{theorem-affine-closure}, that is, when $X=\overline{X^\bullet}^\aff$ is the affine closure of its open orbit (Theorem \ref{thm:heckeIC}), and the monoid $\mathfrak c_X$ is free. This reduction uses the action of the Hecke algebra, and resembles the calculation of the IC sheaf of reductive $L$-monoids in \cite{BNS}, but is much harder. This reduction will be discussed in \S \ref{subsection-reduction-affineclosure} below.

Hence for now let us focus on the case $X=\overline{X^\bullet}^\aff$, assuming that $\mathfrak c_X$ is free. In that case, the maximal possible dimension for the central fiber $\msf Y^{\check\lambda}$ will be called the \emph{critical dimension}, and is equal to 
\[ \frac{1}{2} (\len(\check\lambda)-1),\]
where $\len(\check\lambda) := \sum_i m_i$, for $\check\lambda\in \mathfrak c_X$ written uniquely
as $\sum_i m_i \ch\nu_i$ in terms of our basis of $\mathfrak c_X$. 
Conjecturally, the set of irreducible components of $\msf Y^{\ch\lambda}$ of critical dimension,
ranging over all $\ch\lambda \in \mf c_X$, should correspond to the subset
of the crystal basis 
of $\rho_X$ with weights in $\mf c_X$. 


We do not go quite as far in general, but we show that these irreducible components give rise to the aforementioned \emph{crystal}, in the sense of Kashiwara \cite{Kas93}, over the Langlands dual Lie algebra $\ch{\mathfrak g}$.
Namely, let $\mB_{X^\bullet}^+$ denote the set of irreducible components of 
the central fibers of critical dimension, so $\mB_{X^\bullet}^+$ corresponds to a basis 
of $V_{X^\bullet}^+ := \bigoplus_{\mf c_X} V_{X^\bullet,\ch\lambda}$. 
Formally define $\mB_{X^\bullet}^-$ to be the ``negatives'' of $\mB_{X^\bullet}^+$, so 
$\mB_{X^\bullet}^-$ corresponds to a basis of the dual space $(V_{X^\bullet}^+)^*$. 
Let $\mf B_{X^\bullet} = \mf B_{X^\bullet}^+ \sqcup \mf B_{X^\bullet}^-$. We prove in 
Theorems~\ref{thm:crystal} and \ref{thm:crystal-properties}:

\begin{thm}
\source{intersection-complexes-unramified-l-factors}
\notready 
\label{theorem-crystal}
\source{intersection-complexes-unramified-l-factors}
Let $X=\overline{X^\bullet}^\aff$ satisfy the assumptions of Theorem \ref{theorem-affine-closure}. The set $\mB_{X^\bullet}$ has the structure of a semi-normal, self-dual crystal over 
$\ch{\mf g}$ such that the weights have the properties described in Theorem~\ref{theorem-affine-closure}. 
\end{thm}

We conjecture:

\begin{conj}
\source{intersection-complexes-unramified-l-factors}
\notready 
\label{conjecture-crystal}
\source{intersection-complexes-unramified-l-factors}
The crystal $\mB_{X^\bullet}$ is isomorphic to the unique crystal basis 
of a finite-dimensional $\ch{G}$-module $V_X$. 
\end{conj} 
This would imply Conjecture \ref{conjecture-pushforward-IC}. 

\medskip

Theorem~\ref{theorem-crystal} endows $V_X^+ \oplus (V_X^+)^*$, as we have defined it, 
with an action of an $\SL_2$-triple corresponding to every simple root of $\ch G$. 
These actions imply that the dimensions of the weight spaces 
are invariant under the Weyl group of $G$, which provides a kind of 
``functional equation'' for $\pi_!\Phi_0$. This functional equation can be seen as a geometric analog of the functional equation of the Casselman--Shalika method \cite{C, CS, SaSph}.
The content of 
Conjecture~\ref{conjecture-crystal} is to show that these $\SL_2$-triples satisfy 
the Weyl relations: $[e_\alpha, f_\beta]=0$ for simple roots $\alpha\ne \beta$.
(The Weyl relations imply the Serre relations by \cite[Corollary 4.3.2]{Chriss-Ginzburg}.)   


\smallskip

The construction of the action of the $\SL_2$ corresponding to a simple root $\alpha$ of $G$ 
goes as follows: 
we factor $X \to X \sslash N$ through $X \to X\sslash N_{P_\alpha} \to X\sslash N$, where $P_\alpha$ is the 
sub-minimal parabolic corresponding to $\alpha$ and $N_{P_\alpha}$ is its unipotent radical. 
Then the GIT quotient $X_\alpha:= X\sslash N_{P_\alpha}$ is a spherical variety for the Levi factor $M_\alpha$. 
But now $X_\alpha$ is (usually) larger than the affine closure of its homogeneous part $X_\alpha^\bullet$. 
The irreducible components of $\msf Y_X$ of critical dimension (i.e., elements of $\mB_{X^\bullet}^+$) 
will either go to irreducible components 
\begin{enumerate}
\item of $\msf Y_{X_\alpha^\bullet}$ of critical dimension or 
\item of $\msf Y_{X_\alpha} \sm \msf Y_{X_\alpha^\bullet}$ (not necessarily of critical dimension). 
\end{enumerate}
While the fibers of $\msf Y_X \to \msf Y_{X_\alpha}$ are not necessarily irreducible, we
show that the irreducible components of different relevant fibers can be canonically identified. 
Then we define the $\SL_2$-action by analyzing the two cases above in the base $\msf Y_{X_\alpha}$. 

Under our assumptions, $X_\alpha^\bullet$ is a torus torsor over $\mbb G_m \bs \PGL_2$ and case (i) is an easy
calculation. Meanwhile our study of non-canonical affine embeddings using Hecke actions 
shows that in case (ii) we always get a Mirkovi\'c--Vilonen cycle (i.e., irreducible
component of the intersection of a semi-infinite orbit with a $\msf L^+G$-orbit in the affine Grassmannian). 
The crystal structure on these cycles was constructed by \cite{BGcrys}. 



To check the Weyl/Serre relations, one can similarly reduce to a spherical variety $X^\bullet_{\alpha,\beta}$ for a Levi of semisimple rank two. There are only a handful such varieties (up to center) satisfying our assumptions --- a small subset of the spherical (wonderful) varieties of rank two classified by Wasserman \cite{Wasserman}.  


However, checking the Weyl/Serre relations, even in a few cases, ``by hand'' does not seem to be easy, and we do not have a conceptual proof of them; therefore, we refrain from attempting such a verification. 

The remainder of the introduction will be devoted to describing the two most important elements in the proofs of the theorems above.


\subsection{Reduction to canonical affine closure}
\label{subsection-reduction-affineclosure}
\source{intersection-complexes-unramified-l-factors}

We give an overview of how to reduce the case of an arbitrary affine $X$ with 
$X^\bullet=H\bs G$ to the 
canonical affine closure $X^\can = \ol{H\bs G}^\aff$.  

There is a canonical map $X^\can \to X$, which induces an inclusion 
$X^\can(\mf o) \cap X^\bullet(F) \subset X(\mf o)\cap X^\bullet(F)$ of 
$G(\mf o)$-stable spaces.
Of course, all points of $X^\bullet(F)$ are $G(F)$-translates of points in $X^\can(\mf o)\cap X^\bullet(F)$. 
It is a fact that if $\ch\theta \in \ch\Lambda^-$ is \emph{antidominant} 
and belongs to the monoid $\mf c_X$, then 
the action of the double coset $G(\mf o) t^{\ch\theta} G(\mf o)$ preserves $X(\mf o)\cap X^\bullet(F)$. 
The idea for what follows is that we can obtain 
$X(\mf o)\cap X^\bullet(F)$ by acting on $X^\can(\mf o)\cap X^\bullet(F)$ by $G(\mf o) t^{\ch\theta} G(\mf o)$
for $\ch\theta \in \mf c_X^- := \ch\Lambda^- \cap \mf c_X$.  

The Zastava model $\sY$ lives over $\Bun_B$ and does not carry a Hecke action. 
Thus, to model the $G(F)$-action on $X(\mf o)\cap X^\bullet(F)$ we must use the global model
$\sM=\sM_X$, which lives over $\Bun_G$. 
The canonical map $\sM_{X^\can} \to \sM_X$ is a closed embedding. 
For $\ch\theta \in \mf c_X^-\sm 0$, %(where we remind that $\mf c_X^-$ denotes the antidominant elements of $\mf c_X$)
 let $\sH^{\ch\theta}_{G,C}$ denote the Hecke stack 
over $\Bun_G \times C$ with fibers isomorphic to $\ol\Gr_G^{\ch\theta}$, the closure
of the $\msf L^+G$-orbit in the affine Grassmannian corresponding to $\ch\theta$.  
In reality, we need a symmetrized (multi-point) version of the Hecke stack, but we only describe
the case where there is one point on the curve in this introduction for simplicity. 
There is a well-defined map 
\begin{equation}
 \label{e:Heckeaction}
\source{intersection-complexes-unramified-l-factors}
\sM_{X^\can} \xt_{\Bun_G} \sH_{G,C}^{\ch\theta} \to \sM_X 
\end{equation}
modeling the action of Hecke operators, and we show (Theorem~\ref{thm:comp}) that this map is birational onto its image. 
If we allow multiple points above, then the images of the corresponding Hecke actions, with $\ch\theta$ varying, 
stratify $\sM_X$. 

Under the assumptions of the previous subsection, we show that 
$\sM_{X^\can}$ is irreducible (Corollary~\ref{cor:can-comp}), so the study 
of $\IC_{\sM_X}$ reduces to the study of the Hecke action on $\IC_{\sM_{X^\can}}$ and 
the determination of which of the strata above form irreducible components of $\sM_X$. 
For the latter, we need to understand the closure relations among the different strata (Proposition~\ref{prop:closure-rel}). When $\mf c_{X^\bullet}=\mbb N^r$, the stratum corresponding to $\ch\theta$ 
is contained in the closure of the stratum corresponding to $\ch\theta'$ if and only if $\ch\theta-\ch\theta' \in \mf c_{X^\bullet}$ (more generally, the closure relations are determined by the colors of $X$).  



\subsection{Semi-infinite orbits and dimension estimates}
\label{subsection-dimensions}
\source{intersection-complexes-unramified-l-factors}

There is another way to understand the central fibers $\msf Y^{\check\lambda}$: as subsets of the affine Grassmannian of $G$. 
Let us fix the point $v \in C$ that we take central fibers with respect to. 
Then a $k$-point of $\msf Y^{\ch\lambda}$ is a map $C \to X/B$ such that $C\sm v$ is sent to $X^\circ/B=\pt$. 
Restricting to the completed local ring $\mf o_v$ at $v$ gives a map 
$\msf Y^{\ch\lambda} \to \msf L X^\circ / \msf L^+B$. 
If we fix a base point $x_0\in X^\circ(k)$ to identify $X^\circ \cong B$, we get a map 
$\msf Y^{\ch\lambda} \to \Gr_B$ and this turns out to be a closed embedding. 
The reduced image of the components of $\Gr_B$ in $\Gr_G$ are the
semi-infinite orbits $\msf S^{\ch\lambda}(k) = N(F) t^{\ch\lambda} G(\mf o)/G(\mf o)$. 
After passing to reduced schemes we get identifications 
$\msf Y^{\ch\lambda}_\red = (\msf S^{\ch\lambda} \xt_{\msf L X/ \msf L^+G} \msf L^+ X / \msf L^+G)_\red$ 
and $\olsf Y^{\ch\lambda}_\red = (\olsf S^{\ch\lambda} \xt_{\msf L X /\msf L^+G} \msf L^+ X / \msf L^+ G)_\red$ (see Lemma \ref{lem:Y=ScapGr}). 

\smallskip

Semi-infinite orbits have an important meaning for the geometric Satake equivalence \cite{MV}: the fundamental classes of the irreducible components of the intersection $\msf S^{\check\lambda} \cap \overline\Gr_G^{\check\theta}$, the \emph{Mirkovi\'c--Vilonen cycles}, are in bijection with the ``canonical basis'' for the $\check\lambda$-eigenspace of the irreducible $\ch G$-module of lowest weight $\check\theta$. 

\smallskip

Our analysis of the central fibers $\msf Y^{\ch\lambda}$ is founded upon the following argument from
\cite[\S 3]{MV}. 
The boundary $\olsf S^{\check\lambda}\sm \msf S^{\check\lambda} = \cup_{\check\nu<\check\lambda} \msf S^{\check\nu}$ is a hyperplane section for some projective embedding of $\Gr_G$.
Hence any closed subscheme of $\Gr_G$ which intersects $\olsf S^{\check\lambda}$, also intersects its boundary in codimension one (unless already contained in the boundary). By inductively ``cutting'' by these hyperplanes,
we prove:

\begin{thm}
\source{intersection-complexes-unramified-l-factors}
\notready
\label{theorem-slicing}
\source{intersection-complexes-unramified-l-factors}
Let $X = X^\can$ be as in Theorem \ref{theorem-affine-closure}. 
 Let $\msf b$ be an irreducible component of the central fiber $\msf Y^{\check\lambda}$. Then
 \begin{itemize}
  \item $\dim \msf b \le \frac{1}{2} (\len(\check\lambda)-1)$, 
  \item for a basis element $\ch\nu_i$ of $\mf c_X$ (corresponding to a color), $\msf Y^{\ch\nu_i} = \pt$,
  \item the inequality is an equality only if there is a sequence $\alpha_1,\dots,\alpha_d$ of simple roots (with repetitions) such that $\olsf b \cap \msf S^{\ch\lambda-\check\alpha_1-\cdots-\check\alpha_j}$ is of dimension $\dim \msf b - j$ (hence, also of critical dimension), and $\check\lambda - \sum_{i=1}^d \check\alpha_i = \check\nu$ for a color $\check\nu$.
 \end{itemize}
\end{thm}

The operation of hyperplane ``cutting'' can almost be thought of as the lowering operator 
for some $\SL_2$-triple; unfortunately it is not quite precise enough, see Proposition~\ref{prop:f-hyper}. 

\medskip

If $X \ne X^\can$, then we also show that if $\msf b$ is an irreducible component of 
$\msf Y^{\ch\lambda}_X$ of critical dimension that is not contained in $\msf Y^{\ch\lambda}_{X^\can}$, then
$\ch\lambda$ must be a weight of $V^{\ch\theta_i}$ for one of the $\ch\theta_i$ appearing in 
Theorem~\ref{theorem-general}, and $\msf b$ is birational to a Mirkovi\'c--Vilonen cycle in 
$\msf S^{\ch\lambda} \cap \Gr_G^{\ch\theta_i}$. The latter correspondence comes from the Hecke action \eqref{e:Heckeaction}.

\medskip


Let us comment on how the above relates to Theorems \ref{theorem-semismall} and \ref{theorem-general}. 
Under our assumption that $\mf c_{X^\bullet} = \mbb N^r$, the space $\sY^{\ch\lambda}_{X^\can}$ 
is irreducible. 
Then the dimension estimate in Theorem~\ref{theorem-slicing}, together with a factorization property 
of $\sY_X$, implies the ``stratified semi-smallness'' condition. 

The irreducible components of $\msf Y_{X^\can}$ of critical dimension 
correspond to a basis of $V^+_{X^\bullet}$. The other irreducible components of $\msf Y_X \sm \msf Y_{X^\can}$
of critical dimension correspond to Mirkovi\'c--Vilonen cycles in $\Gr_G^{\ch\theta_i}$, and
these provide a basis for $V^{\ch\theta_i}$ in Theorem~\ref{theorem-general} by geometric Satake. 

\subsection{Organization of the paper}
In \S\ref{section-spherical} we briskly review the salient combinatorics
of spherical varieties and the classification of $G(\mf o)$-orbits of the loop
space of $X^\bullet$. In \S\ref{sect:models} we introduce the
global and Zastava models for the arc space of $X$ and their stratifications, explain why they are indeed models,
and prove some foundational properties. In \S\ref{sect:compactification}, we
introduce the compactification of the Zastava model and define the
central fibers of (compactified and non-compactified) Zastava models. Then we 
perform the comparison between $\pi_!\IC_{\sY}$ and $\bar\pi_! \IC_{\barY}$ 
that accounts for the ``numerator'' in the Euler factor \eqref{equation-pushforward-IC-expl}. 

Sections \ref{sect:Heckeact} and \ref{sect:centralfiber} are the 
technical heart of this paper. In \S\ref{sect:Heckeact} we establish 
the closure relations for the global model $\sM_X$ and determine its irreducible
components. This involves a study of the $G$-Hecke action on the global model,
which also reduces the problem to the canonical affine closure, as explained 
earlier. In \S\ref{sect:centralfiber}, we analyze the geometry of the
central fiber and prove the crucial dimension estimates using the
Mirkovi\'c--Vilonen boundary hyperplanes of semi-infinite orbits. 
This allows us to prove Theorem~\ref{theorem-semismall}. 

In \S\ref{sect:crystal}, we prove the aforementioned results on crystals. 
In \S\ref{sect:nearby}, we combine the results of the preceding sections to 
compute the nearby cycles of the IC complex on the global model using a 
well-known contraction principle. Here we establish that 
the nearby cycles functor does indeed correspond to the asymptotics map under 
the sheaf--function dictionary. 

In Appendix \ref{appendix:A}, we collect various technical results 
concerning the stratification of the global model, some of which use the
notion of generic-Hecke modification from \cite{GN}. 


\subsection{Index of notation}\  

In general, we will use calligraphic 
letters $\Cal D, \Cal V$ to denote standard combinatorial objects associated
to spherical varieties in the literature, script letters $\Scr M, \Scr Y, \Scr F$ for
algebraic stacks and sheaves, sans serif letters $\msf Y, \msf S$ for (ind-)schemes 
that are subspaces of certain loop spaces with respect to a fixed point $v\in \abs C$. The following table contains most of the notation used in this paper, except for notation defined and used locally.


\begin{longtable}{>{\raggedleft}p{.18\textwidth} p{.78\textwidth} }
$k$ & an algebraically closed field. The characteristic of $k$ can be zero or positive, but in the latter case we will impose some restrictions on our spherical varieties (see \S \ref{assumptions-finitefield}), to ensure that their geometry is similar to that in characteristic zero. \\
{$\mathbb F$, $\Fr$} & {At some points in this paper, $k$ is the algebraic closure of a finite field $\mathbb F$, and then $\Fr$ denotes the geometric Frobenius morphism}. \\
$\pt$ & $\Spec k$. \\ 
{$C$} & {a connected smooth projective curve over $k$.} \\ 
{$\Sym C,\, C^{(n)}$} & {the scheme of effective divisors on (=symmetric powers of) $C$, and the component of divisors of degree $n$.} \\
{$\oo C^n, \oo C^{(n)}$} & {the open subsets of distinct $n$-tuples of points, resp.\ multiplicity free divisors of degree $n$, on the curve.} \\
{$\oo \prod$, $\oo\times$} & {for schemes living over any partially symmetrized powers of the curve, the restriction of their Cartesian product over the multiplicity-free locus.} \\
{$\Bbbk = k(C)$} & {the field of rational functions on $C$.}  \\
{$\abs C = C(k)$} & the set of closed points of $C$. \\
{$\mf o_v$} & {for $v\in \abs C$, it denotes the completion of the local ring at $v$.} \\
{$F_v$} & {the fraction field of $\mf o_v$.
By choosing a local coordinate $t$ we have a non-canonical isomorphism 
$\mf o_v \cong k\tbrac t =: \mf o$ and $F_v \cong k\lbrac t =: F$. We sometimes implicitly
make this identification when the choice of local coordinate is irrelevant.} \\
{$\mbb N$} & {the monoid of non-negative integers.} \\
& \\ 
{$G$} & {a connected reductive group over $k$.} \\
{$T$} & {the (abstract) Cartan of $G$, i.e., the reductive quotient of any Borel subgroup. We sometimes fix a splitting $T\hookrightarrow B\hookrightarrow G$ of the abstract Cartan into a Borel subgroup.} \\
{$\Cal B$} & {the flag variety of Borel subgroups of $G$.} \\
{$W$} & {the (abstract) Weyl group of $G$.} \\
{$s_\alpha \in W$} & {for a simple root $\alpha$, the corresponding reflection.} \\
{$\check \Lambda_G$ ($\Lambda_G$)} & {the coweight (resp.~weight) lattice of $T$. The index $G$ will often be omitted.} \\
{$\check\Lambda^+_G$ ($\Lambda_G^+,\, \ch\Lambda^-_G$)} & 
The monoid of dominant coweights (resp., dominant weights, antidominant coweights). \\
{$\ch\Lambda^\pos_G$ ($\Lambda^\pos_G \subset \Lambda_G$)} & 
The monoid generated by the non-negative \emph{integral} span of the positive coroots
(resp.~roots) in $\ch\Lambda_G$.  \\
{$\check\Delta_G$ ($\Delta_G$)} & {the set of simple coroots (resp.~roots)
of $G$.} \\
{$2\ch\rho_G \in \ch\Lambda_G$ ($2\rho_G \in \Lambda_G$)} & {the sum of 
the positive coroots (roots) of $G$} \\
{$\ch\lambda \ge \ch\mu$} & {For $\ch\lambda,\ch\mu \in \ch\Lambda_G$, this means that $\ch\lambda - \ch\mu \in \ch\Lambda^\pos_G$.} \\
{$\ch G$} & {the Langlands dual group of $G$ over $\ol\bbQ_\ell$, i.e., $\ch G$ is
the connected reductive group where the weights, roots of $\ch G$ equal the 
coweights, coroots of $G$, etc.} \\
{$V^{\ch\lambda}$} & {for $\ch\lambda \in \ch\Lambda_G$ either dominant or antidominant, this denotes
the irreducible $\ch G$-module over $\ol\bbQ_\ell$
with highest (resp.~lowest) weight $\ch\lambda$.
Similarly, $V^\lambda$ denotes the irreducible $G$-module over $k$ for $\lambda\in \Lambda_G^+$.}
\\
& \\
{variety} & {will mean a reduced, finite type $k$-scheme (not necessarily irreducible).}  \\
{$\Scr Z/H$} & {for a stack $\Scr Z$ with an action of an algebraic group $H$, this will denote the quotient stack.} \\
{$Z\sslash H$} & {for an affine variety $Z$ over $k$, and a group $H$ acting on it, the invariant-theoretic quotient $\Spec k[Z]^H$.} \\ 
{$\Scr X \xt^G \Scr Y$} & 
if $\Scr X, \Scr Y$ are stacks with (right) $G$-actions, we will use  this 
to denote the stack quotient $(\Scr X \xt \Scr Y)/G$ by the diagonal action. \\
{$\msf L^+ X,\, \msf L X$} & {the formal arc and loop spaces of a scheme $X$ (see \S \ref{sect:formalloop}).} 
\\
& \\
{$\rmD^b_c(\sZ)$} & {for an algebraic stack $\sZ$, this is the derived category of 
bounded constructible $\ol\bbQ_\ell$-complexes on $\sZ$.} \\
{`sheaf'} & {means a complex of sheaves. 
All functors between sheaves are derived functors.} \\
{$\rmP(\sZ) \subset \rmD^b_c(\sZ)$} & {when $\sZ$ is locally of finite type over $k$, this is the abelian category of perverse sheaves.}\\
{$^p\rmD^{\le 0}, {^p}\rmD^{\ge 0}$} & 
the subcategories with respect
to the perverse $t$-structure. \\
{$\IC_{\Scr Z}$} & {the direct sum of the intersection cohomology complexes
of all irreducible components of $\Scr Z$. When working over a finite field, we will normalize this sheaf to be pure of weight zero.} 
\\
&\\
{$X$} & {an affine $G$-spherical variety over $k$. (See \S \ref{sect:spherical} for notions pertaining to spherical varieties.)} \\
{$X^\circ$} & {the open $B$-orbit, for a fixed choice of Borel subgroup $B$.} \\
{$x_0 \in X^\circ(k)$} & {a fixed base point.} \\
{$H$} & {the stabilizer of $x_0$.} \\
{$X^\bullet$} & {the open $G$-orbit $H\bs G $.} \\
{$X^\can$} & {the ``canonical'' affine embedding $\Spec k[X^\bullet]$ of $X^\bullet$.} \\
{$T_X$} & {the (abstract) Cartan of $X$, that is, the quotient by which the abstact Cartan of $G$ acts on $X^\circ\sslash N$, where $N$ is the unipotent radical of $B$.} \\
{$\Lambda_X, \ch \Lambda_X$} & {the character and cocharacter groups of $T_X$. Our assumptions on $X$ will identify $\ch\Lambda_X$ with $\ch\Lambda_G$, so it will often just be denoted by $\ch\Lambda$.} \\
{$\check G_X$} & {the dual group of $X$; it has a canonical maximal torus isomorphic to the dual of $T_X$.} \\
{$\Cal V$} & {the cone of invariant valuations of $X$; equivalently, the antidominant chamber of the dual group of $X$.} \\
{$\mathfrak c_X^\vee \subset \Lambda_X$ ($\mathfrak c_X\subset \ch\Lambda_X$)} & 
{the monoid of weights of $T_X$ on $k[X\sslash N]$ (resp., its dual monoid).} \\
{$\Cal C_0=\Cal C_0(X)\subset \mathfrak t_X$} & 
{the cone spanned by $\mathfrak c_X$, inside of the vector space spanned by $\ch\Lambda_X$.} \\
{$\mathfrak c_X^-$} & {the intersection of $\mathfrak c_X$ with the cone $\Cal V$ of invariant valuations.} \\
{$\mathcal D(X)$} & {the set of irreducible $B$-stable divisors in $X$.} \\
{$\mathcal D$} & {the set of colors, i.e., irreducible $B$-stable divisors which are not $G$-stable; equivalently, this can be identified with $\mathcal D(X^\bullet)$.} \\
{$\mathcal D(\alpha)$} & {for a simple root $\alpha$, the set of colors $D$ of $X^\bullet$ such that $D P_\alpha \supset X^\circ$, 
where $P_\alpha$ is the parabolic generated by $B$ and the root space $\mathfrak g_{-\alpha}$.}\\
{$\varrho_X(D)=\ch\nu_D$}&{for $D\in \mathcal D(X)$, the associated $B$-invariant valuation, restricted to the group of nonzero $B$-eigenfunctions: $\check \nu_D: k(X)^{(B)} \to \mathbb Z$, and understood as  a functional on the character group $\Lambda_X = k(X)^{(B)}/k^\times$.}\\
{$\mathfrak c_X^{\Cal D} \subset\ch\Lambda_X$} & 
the monoid generated by the $\ch\nu_D,\, D\in \Cal D$.  \\
{$\ch\lambda\preceq \ch\mu$} & {for $\ch\lambda,\ch\mu\in \ch\Lambda_X$, this means that 
$\ch\mu-\ch\lambda \in \mathfrak c_X^{\Cal D}$.} \\
{$\Cal D^G_{\mathrm{sat}}(X)$} & {the set of those primitive (=indecomposable) elements in $\mathfrak c_X^-$ that
are minimal with respect to the $\preceq$ partial order.}
\\
&\\
{$\Bun_G, \Bun_B$} & {the moduli stack of $G$-bundles, resp.\ $B$-bundles, on $C$.} \\
{$\sM_{X}$} & the ``global model'' of generic maps from a curve to $X/G$; it lives over $\Bun_G$. The restriction of such a map, defined over $k$, to the formal neighborhood of a point $v\in |C|$ gives rise to a well-defined ``valuation'', that is, an element of $(X^\bullet(F_v)\cap X(\mf o_v))/G(\mf o_v)$. 
\\
& See Section \ref{sect:models} for the various models of the arc space. For any model, when $X$ is understood, the index will be omitted. \\
{$\sY_X$} & {the ``Zastava model'' of generic maps from a curve to $X/B$; it lives over $\Bun_B$. The restriction of such a map, defined over $k$, to the formal neighborhood of a point $v\in |C|$ gives rise to a well-defined element of $(X^\circ(F_v)\cap X(\mf o_v))/B(\mf o_v)$} \\
{$\ol\Bun_B$} & {Drinfeld's compactification of $\Bun_B$, see \S\ref{sect:basicproperties-comp}.}
\\
{$\barY_X$}& {the compactified Zastava model (Section \ref{sect:compactification}).}
\\
{$\Scr A$} & {the global/Zastava model for the $T_X$-space $X\sslash N$.}
\\
{$\pi,\bar\pi$} & the natural maps $\pi:\sY_X\to\Scr A$, $\bar\pi:\barY_X\to\Scr A$ (extending $\pi$). 
\\
$X^\bullet(F)_{G:\ch\theta}$, $\msf L^{\ch\theta}X$ & 
{the $G(\mf o)$-orbit on $X^\bullet(F)$ parametrized by $\ch\theta\in \Cal V \cap \ch\Lambda_X$ (Theorem \ref{thm:Gorbits}), and the corresponding stratum of the loop space. When $\ch\theta \in \mathfrak c_X^-$, these belong to $X(\mf o)$, resp.\ the arc space $\msf L^+X$.}
\\
{$\Sym^\infty(\msf S)$} & 
{the set of multisets in elements of a set $\msf S$; equivalently, the free monoid $\bigoplus_{\msf S}\mathbb N$ in the elements of $\msf S$.}
\\
{$C^{\mf P}, \oo C^{\mf P}$} & 
{for a multiset $\mf P = \sum_{s\in \msf S} N_s [s]$, the partially symmetrized power $\prod_{\msf S} C^{(N_s)}$ of the curve, and its disjoint locus $\oo\prod_{\msf S} \oo C^{(N_s)}$.}
\\
{$\sM^{\ch\Theta}$} & 
{for $\ch\Theta \in \Sym^\infty(\mathfrak c_X^-\sm \{0\})$, the stratum of $\sM$ containing those maps whose multiset of nontrivial valuations (as elements of $\mathfrak c_X^-\sm \{0\}$) is equal to $\ch\Theta$. When $\ch\Theta=\{\ch\theta\}$ is a singleton, we will write $\sM^{\ch\theta}$.}
\\
{$\Scr A^{\ch\lambda}$} & {the connected component of $\Scr A$ of maps with total valuation $\ch\lambda \in \mathfrak c_X \subset \ch\Lambda_X= T_X(F)/T_X(\mf o)$.} 
\\
{$\sY^{\ch\lambda}, \barY^{\ch\lambda}$} & {the preimage of $\Scr A^{\ch\lambda}$ in $\sY$, resp.\ in $\barY$. They live over strata $\Bun_B^{-\ch\lambda}$, $\ol\Bun_B^{-\ch\lambda}$ of $\Bun_B$, resp.\ $\ol\Bun_B$.}
\\
{$\sY^{\ch\lambda,\ch\Theta},\, \barY^{\ch\lambda,\ch\Theta}$} & {the fiber products of $\sY^{\ch\lambda}, \barY^{\ch\lambda}$ with $\sM^{\ch\Theta}$ over $\sM$.}
\\
{$\sY^D$} & {for $D \in \mathbb N^{\mathcal D}$, a certain connected/irreducible component of the ``open Zastava'' space $\sY_{X^\bullet} = \sY^{?,0}$, defined in \S \ref{sect:Y0cc}. The question mark $?$ corresponds to the valuation $\varrho_X(D)$.}
\\
{$\barM^{\ch\Theta}$} & {denotes the closure of the stratum $\sM^{\ch\Theta}$. Note that $\barY^{\ch\lambda}$, $\barY^{\ch\lambda,\ch\Theta}$, in contrast, are \emph{not} closures of strata, but strata of the compactified Zastava space. In the case of the global model, there is no room for confusion, so we allow ourselves this notation, for typographical reasons.}
\\
$C_{\ch\nu}$, ${_{\ch\nu}}\ol\Bun_B^{\ch\lambda}$, ${}_{\ch\nu} \barY^{\ch\lambda}$ & 
for $\ch\nu\in \ch\Lambda_G^{\pos}$, the partially symmetrized power $C^{\mf P}$ when $\ch\nu$ is thought of as a multiset $\mf P$ in the simple coroots, a stratum of $\ol\Bun_B^{\ch\lambda}$, and a stratum of $\barY^{\ch\lambda}$, isomorphic, respectively, to $C_{\ch\nu} \times \Bun^{\ch\lambda+\ch\nu}_B \into \ol\Bun^{\ch\lambda}_B$ and  $C_{\ch\nu} \times \Scr Y^{\ch\lambda-\ch\nu}$ (see \S \ref{sect:strat-compactified}).
\\
& \\
{$\Gr_G, \Gr_B$} & the affine Grassmannian of $G$, resp.\ $B$. 
\\
$\Gr_{G,\Sym C}$, $\Gr_{B,\Sym C}$ & {the Beilinson--Drinfeld affine Grassmannians, living over $\Sym C$ (see \S \ref{sect:YGr}).}
\\
{$\Gr_G^{\ch\theta}$, $\ol\Gr_G^{\ch\theta}$} & {for $\ch\theta\in \ch\Lambda_G^-$, the $\msf L^+G$-orbit in the affine Grassmannian containing the class of $t^{\ch\theta}$, and its closure.} \\
{$\ol\Gr^{\ch\Theta}_{G,C^{\ch\Theta}}$} & {for $\ch\Theta \in \Sym^\infty(\ch\Lambda_G^- \sm 0)$, the multi-point version of $\ol\Gr^{\ch\theta}$, see \S \ref{sect:strata-Grassmannian}.} \\
{$\msf S^{\ch\lambda}$, $\olsf S^{\ch\lambda}$} & {the ``semi-infinite'' $\msf LN$-orbit of $t^{\ch\lambda}$ in $\Gr_G$, and its closure.}
\\
{$\msf Y^{\ch\lambda},\, \olsf Y^{\ch\lambda}$} & {the central fibers of $\sY^{\ch\lambda}, \barY^{\ch\lambda}$, living over a point of the diagonal stratum $C\hookrightarrow \Scr A^{\ch\lambda}$. They can be identified as subspaces of $\msf S^{\ch\lambda}$, $\olsf S^{\ch\lambda}$ (see \S \ref{def:centralfiber}). From \S\ref{def:centralfiber} onwards, we will implicitly consider
only the underlying reduced structure on all central fibers.
}
\\
& \\
{$\mB_{X,\ch\lambda}$} & {for $\ch\lambda\in \mathfrak c_X$, the set of all irreducible components of critical dimension of the central fiber $\msf Y^{\ch\lambda}$, see Proposition \ref{prop:Vbasis}.} \\
{$\mB_X^+, \mB_X$} & {the union of all $\mB_{X,\ch\lambda}$, $\ch\lambda\in \mf c_X$, and the ``crystal of $X$'' (\S \ref{sect:defcrys}).}
\end{longtable}





\subsection{Acknowledgments}

We thank R.~Bezrukavnikov, V.~Drinfeld, T.~Feng, M.~Finkelberg, V.~Ginzburg, F.~Knop and D.~Gaitsgory for helpful comments and conversations. We thank the anonymous referee for helpful suggestions. We thank the Institute for Advanced Study for its hospitality during the academic year 2017--2018, during which part of our work was conducted. Y.S.~was supported by NSF grants DMS-1801429 and DMS-1939672, and by a stipend to the IAS from the Charles Simonyi Endowment.
J.W.~was supported by NSF grant DMS-1803173.
